\section{The response system and the quartic kernel}
\label{app:response}

\subsection{The condensate as a source}
\label{app:response-source}

Just above the onset, $W_x = \eta\,
w_0(r)\,\psi(x,y)$ with $W_y = -iW_x$, and $|\psi|^2 = \sum_G \omega_G\,
e^{iG\cdot x}$ with $\omega_G = (-1)^{m+n+mn}e^{-\gamma_G/4}$ on the reciprocal
lattice, $G = G_{mn}$, so that the quartic free energy carries the weights
$|\omega_G|^2 = e^{-\gamma_G/2}$.
At second order in $\eta$ the condensate sources the abelian field and
the metric at each $G$, through the current $\epsilon^{3bc}W^b\,DW^c$ and the
stress tensor of $W$ evaluated on the zero mode. Both close on
$|\psi|^2$ alone, so the response to the harmonic $\eta^2\omega_G\,e^{iGx}$, which we take along $x$, is
\begin{equation}
\begin{split}
h_{MN}\,dx^Mdx^N &= \eta^2\omega_G\,e^{iGx}\Big[-U H_{tt}\,dt^2 + e^{2V}\big(H_{xx}\,dx^2 + H_{yy}\,dy^2\big) + e^{2Z}H_{zz}\,dz^2 \\
&\qquad\qquad\qquad + \frac{H_{rr}}{U}\,dr^2 + 2ie^{2V}H_{xr}\,dx\,dr\Big]~, \\
\delta A^3 &= \eta^2\omega_G\,e^{iGx}\,a_y\,dy~, \qquad \hat b = iG\,a_y~,
\end{split}
\end{equation}
with $\hat b$ the induced field strength. Per unit $\eta^2\omega_G e^{iGx}$
the stress tensor of $W$, in the normalisation of the Einstein equations
$G_{MN} - 6g_{MN} = \alpha\,(F_{MP}F_N{}^P - \tfrac14 g_{MN}F^2)$, is diagonal,
%% check: source
\begin{equation}
\begin{split}
&S_{tt} = Ue^{-4V}\big(Ue^{2V}w_0'^2 - Bw_0^2\big)~, \qquad
S_{xx} = S_{yy} = -Be^{-2V}w_0^2~, \qquad \\
&S_{zz} = e^{2Z-4V}\big(Bw_0^2 - Ue^{2V}w_0'^2\big)~, \qquad
S_{rr} = \frac{e^{-4V}}{U}\big(Bw_0^2 + Ue^{2V}w_0'^2\big)~,
\end{split}
\end{equation}
and it is conserved on the zero-mode equation, $\nabla^M S_{MN} = 0$ up to
the Lorentz force on the condensate current, whose source for $\hat b$ is
$s\,p_1w_0^2$, $s = G^2$. Since the background depends on $r$ only, the response at
wavenumber $G$ is a radial boundary value problem in which $G$ enters
through $s$ alone, and the quartic free energy is the on-shell action of the
response paired against its own source.

\subsection{The abelian sector}
\label{app:response-abelian}

With $p_1 = e^{Z-2V}$ and $p_2 = Ue^Z$, the
induced abelian field $\hat b$ at harmonic $s$ obeys, with the metric frozen,
\begin{equation}
\big(p_2\,\hat b'\big)' = s\,p_1\big(\hat b - w_0^2\big)~,
\end{equation}
regular at the horizon and vanishing at the boundary, and
$X_{\rm ab}(s) = \int p_1 w_0^2\,\hat b\,dr$. Writing $\mathcal{A} = -\tfrac{d}{dr}p_2\tfrac{d}{dr}
\ge 0$, $\mathcal{M} = p_1$, $\mathcal{T} = \mathcal{M}^{-1/2}\mathcal{A}\mathcal{M}^{-1/2}$ and $\mathbf{u} = \mathcal{M}^{1/2}w_0^2$, this is
$(C_4 - X_{\rm ab})(s) = \langle \mathbf{u}, \mathcal{T}(\mathcal{T}+s)^{-1}\mathbf{u}\rangle$, a Stieltjes transform of a
positive measure, hence completely monotone in $s$ on every brane. The
identity $\tfrac{d}{dr}\big[p_2 w^3 w'\big] = 3p_2 w^2 w'^2 - Bp_1 w^4$ on
solutions of the zero-mode equation, whose endpoint terms vanish, gives
$\langle \mathbf{u}, \mathcal{T}\mathbf{u}\rangle = \int p_2\big((w_0^2)'\big)^2 dr = \tfrac43 B_c C_4$
and hence the tail $(C_4 - X_{\rm ab})(s) = \tfrac43 B_c C_4/s + O(s^{-2})$
on every brane.

\subsection{The coupled metric and photon response}
\label{app:response-metric}

The classifying symmetry of the
linearised system is $y$-parity combined with the flip of the $U(1)$ charge,
since the background $F_{xy}$ is odd under $y$-parity alone. Under it $a_y$
is even and $a_x$, $a_r$ are odd and unsourced. The diffeomorphisms
$\xi = (\xi^x, \xi^r)\,e^{iGx}$ of the even sector shift $H_{tt}$ by a
multiple of $\xi^r$ and $H_{xx}$ by multiples of $\xi^x$ and $\xi^r$, with
no derivatives, so the algebraic gauge $H_{tt} = H_{xx} = 0$ can be reached
pointwise. The linearised equations are
$\delta(G_{MN} - 6g_{MN}) - \alpha\,\delta T_{MN} = \alpha S_{MN}$, with
$\delta T$ the variation of the abelian stress in both the metric and the
field, and the $y$ component of the linearised Maxwell equation sourced by
$J$. Using the background constraint
$U'(2V' + Z') + 2UV'(V' + 2Z') = 12 - \alpha B^2e^{-4V}$ to simplify the
coefficients, the $(rr)$ and $(xr)$ components are first order,
%% check: rr
\begin{equation}
\begin{split}
&\big(\tfrac12\alpha B^2e^{-4V} - 6\big)H_{rr} + \tfrac12 G\,(U' + 2UV' + 2UZ')\,H_{xr}
+ \tfrac14(U' + 2UV' + 2UZ')\,H_{yy}' \\
&\quad + \tfrac14(U' + 4UV')\,H_{zz}' - \tfrac12\big(\alpha B^2e^{-4V} + G^2e^{-2V}\big)H_{yy} - \tfrac12 G^2e^{-2V}H_{zz} \\
&\quad + \alpha Be^{-4V}\,\hat b = \alpha e^{-4V}\big(Bw_0^2 + Ue^{2V}w_0'^2\big)~,
\end{split}
\end{equation}
%% check: xr
\begin{equation}
\begin{split}
&\tfrac14 G^2(U' + 2UV' + 2UZ')\,H_{rr} + \alpha B^2 G\,Ue^{-2V}H_{xr}
- \tfrac12 G^2U\big(H_{yy}' + H_{zz}'\big) \\
&\quad + \tfrac12 G^2U(V' - Z')\,H_{zz} + \alpha BUe^{-2V}\,\hat b' = 0~,
\end{split}
\end{equation}
the $(yy)$ and $(zz)$ components are second order,
%% check: yy
\begin{equation}
\begin{split}
&\tfrac12 U H_{zz}'' + \tfrac12(U' + UV' + 2UZ')\,H_{zz}' - \tfrac12 G^2e^{-2V}H_{zz}
+ \tfrac12\alpha B^2e^{-4V}H_{yy} + GUH_{xr}' \\
&\quad + G(U' + 2UV' + UZ')\,H_{xr} - \tfrac14(U' + 2UV' + 2UZ')\,H_{rr}' \\
&\quad - \tfrac12\big(G^2e^{-2V} + \alpha B^2e^{-4V} + 12\big)H_{rr}
- \alpha Be^{-4V}\,\hat b = -\alpha Be^{-4V}w_0^2~,
\end{split}
\end{equation}
%% check: zz
\begin{equation}
\begin{split}
&\tfrac12 U H_{yy}'' + \tfrac12(U' + 3UV')\,H_{yy}' - \tfrac12\big(\alpha B^2e^{-4V} + G^2e^{-2V}\big)H_{yy}
+ GUH_{xr}' \\
&\quad + G(U' + 3UV')\,H_{xr} - \tfrac14(U' + 4UV')\,H_{rr}' + \tfrac12\big(\alpha B^2e^{-4V} - G^2e^{-2V} - 12\big)H_{rr} \\
&\quad + \alpha Be^{-4V}\,\hat b = \alpha e^{-4V}\big(Bw_0^2 - Ue^{2V}w_0'^2\big)~,
\end{split}
\end{equation}
and the Maxwell equation is
%% check: My
\begin{equation}
\big(p_2\,\hat b'\big)' - s\,p_1\big(\hat b - w_0^2\big) + BG\big(p_2 H_{xr}\big)'
+ \tfrac12 sBp_1\big(H_{yy} - H_{zz} - H_{rr}\big) = 0~,
\end{equation}
which with the metric perturbation switched off is the abelian equation of appendix~\ref{app:response-abelian}.
These five rows, with the derivatives of the two first-order ones, form a
linear system for $(H_{rr}, H_{xr}, H_{rr}', H_{xr}', H_{yy}'', H_{zz}'', \hat b'')$
that is non-singular on the brane. The even sector therefore closes on the
three evolution fields $(H_{yy}, H_{zz}, \hat b)$, with $H_{rr}$ and $H_{xr}$
fixed algebraically; the $(tt)$ and $(xx)$ equations then vanish identically
on the solved forms, as the linearised Bianchi identity and source
conservation require. Regularity at the horizon fixes the slope of
each evolution field in terms of its value, and the regular solutions have
$H_{rr}(r_p) = 0$, so that $h_{rr} = H_{rr}/U$ is regular and the
temperature is not modulated. The boundary indicial exponents are
$\{0, 4\}$ for the metric fields and $\{0, 2\}$ for $\hat b$; the
normalisable solution is the one with the three leading coefficients zero,
and then $H_{yy}, H_{zz}, H_{rr} = O(r^{-4})$, $H_{xr} = O(r^{-5})$ and
$\hat b = O(r^{-2})$. At $\beta = 0$ the $(H_{yy}, H_{zz})$ block reduces to
the probe graviton system and the Maxwell row to the abelian equation.

\paragraph{The pairing.} Let $\mathcal{L}_q$ be the part of
$\sqrt{-g}\,\big(R + 12 - \tfrac{\alpha}{2}F^2\big)$ quadratic in the
response $(h, a_y)$ (the harmonic times its conjugate, averaged over
the cell), and add the couplings to the condensate,
\begin{equation}
\mathcal{L}_2 = \mathcal{L}_q + \alpha\sqrt{-g}\,S^{MN}\big(h_{MN} + \bar h_{MN}\big)
+ 2\alpha\,p_1w_0^2\big(\hat b + \bar{\hat b}\big)~,
\end{equation}
with indices raised by the background metric and $h_{MN}$ stripped of
$\eta^2\omega_Ge^{iGx}$. The coefficients of the two source terms are not
chosen: they are the unique ones for which the Euler--Lagrange equations of
$\mathcal{L}_2$ are the five rows above, with a single overall factor
$-\sqrt{-g}$ on the Einstein rows and $-2\alpha$ on the Maxwell row.
Since $\mathcal{L}_q$ is a homogeneous quadratic form, Euler's theorem gives
$\sum_\phi \phi\,\mathrm{EL}_\phi(\mathcal{L}_q) = 2\mathcal{L}_q - \mathcal{Q}'$, the sum running over the fields and their conjugates, with
\begin{equation}
\mathcal{Q} = \sum_\phi\Big[\phi\Big(\frac{\partial\mathcal{L}_q}{\partial\phi'}
- \Big(\frac{\partial\mathcal{L}_q}{\partial\phi''}\Big)'\Big)
+ \phi'\,\frac{\partial\mathcal{L}_q}{\partial\phi''}\Big]~,
\end{equation}
so that on shell $\mathcal{L}_2$ is half its source terms plus
$\tfrac12\mathcal{Q}'$. The flux vanishes at both ends. At the horizon every
term of $\mathcal{Q}$ carries a factor $U$ except
%% check: Qhor
\begin{equation}
\begin{split}
\mathcal{Q}(r_p) = -\tfrac32\,U'e^{2V+Z}\Big[&H_{rr}\bar H_{rr} - \tfrac16 H_{rr}\big(\bar H_{yy} + \bar H_{zz}\big) \\
&- \tfrac16\bar H_{rr}\big(H_{yy} + H_{zz}\big)\Big]~,
\end{split}
\end{equation}
which vanishes with $H_{rr}(r_p)$; at the boundary it falls off as $r^{-2}$,
set by the photon. With the fall-offs above every local boundary term
quadratic in the response, the Gibbons--Hawking term and the holographic
counterterms including the logarithmic $F^2$ one, vanishes as the cutoff is
removed, and the horizon is not a boundary, so the quartic free energy is
the bulk on-shell action alone. In the units of the probe kernel,
where the $\beta \to 0$ limit is $C_4 - X_{\rm ab}$, the on-shell action is
multiplied by $-1/2\alpha$, and the source terms give
\begin{equation}
K(s;\beta) = C_4 - X(s;\beta) - \tfrac12\,\mathcal{P}(s;\beta)~,
\end{equation}
with
\begin{equation}
X = \mathrm{Re}\!\int p_1w_0^2\,\hat b\,dr~, \qquad
\mathcal{P} = \mathrm{Re}\!\int\sqrt{-g}\;S^{MN}\,\bar h_{MN}\,dr~,
\end{equation}
where $\hat b$ and $h$ are the coupled solution, including the algebraic
components, and the contact term $C_4$ is the quartic self-interaction of
$W$. The weight $\tfrac12$ on $\mathcal{P}$ relative to $X$ is the
ratio of the two source couplings in $\mathcal{L}_2$, which the field
equations fix. The difference between $X$ and the abelian-only $X_{\rm ab}$
of appendix~\ref{app:response-abelian} is the part of the photon response
induced through the metric, and it is $X_{\rm ab}$ that carries the spectral
representation and the tail identity.

\paragraph{The weight $\tfrac12$.} The weight has the familiar origin of
the matter couplings: a stress tensor couples to the metric through
$\tfrac12\sqrt{-g}\,T^{MN}h_{MN}$, with
$T^{MN} = (2/\sqrt{-g})\,\delta S/\delta g_{MN} = 2\alpha S^{MN}$ in the
normalisation above, while a current couples to the gauge field through
$\sqrt{-g}\,J^Ma_M$ with no such factor. Gauge invariance fixes it
independently of the action. Under $\xi = (\xi^x, \xi^r)\,e^{iGx}$ the response
shifts by $h \to h + \mathcal{L}_\xi g$ and $a_y \to a_y + B\xi^x$. For a
weight $w$ on $\mathcal{P}$, the shift of $X + w\mathcal{P}$ by $\xi^x$ is
proportional to $w - \tfrac12$ pointwise, because the $(xx)$ component of
the condensate stress is the magnetic pressure of the current that couples
to $a_y$. At $w = \tfrac12$ the whole shift is a total derivative,
%% check: Qxi
\begin{equation}
\delta_\xi\Big(X + \tfrac12\mathcal{P}\Big) = \Big[\mathcal{Q}_\xi\Big]_{r_p}^{\infty}~, \qquad
\mathcal{Q}_\xi = e^{Z-2V}\big(Bw_0^2 + Ue^{2V}w_0'^2\big)\,\mathrm{Re}\,\xi^r~,
\end{equation}
which vanishes for every $\xi$ that keeps the horizon in place,
$\xi^r(r_p) = 0$ as fixed temperature requires, and grows more slowly than
$r^3$ at the boundary, which includes the boundary Weyl rescalings. The
kernel is therefore the same in every such gauge. Transforming the solutions to
radial gauge, $H_{rr} = H_{xr} = 0$, and to a generic smooth gauge confirms
this (table~\ref{tab:checks}).

Pairing the metric response instead with the full stress, including the
Maxwell stress $T_{MN}[\hat b_{\rm ab}]$ of the decoupled photon, gives the
dressed pairing $\mathcal{P}_{\rm dress}$, with
$X - X_{\rm ab} = \tfrac12\big(\mathcal{P}_{\rm dress} - \mathcal{P}\big)$,
because the cross term of $\mathcal{L}_q$ between $h$ and $a_y$ is $\alpha$
times the dressing pairing up to a total derivative. Hence
$K = C_4 - X_{\rm ab} - \tfrac12\mathcal{P}_{\rm dress}$, a split into
gauge-invariant parts, which agrees with $K$ only at weight $\tfrac12$.

The uniform response $K_{G=0}$ of section~\ref{sec:5.1} comes from a
separately derived $G = 0$ system in the gauge $H_{tt} = 0$, as a
solvability condition: it is the contact term plus the shift of the onset
eigenvalue by the condensate's own uniform metric response,
$K_{G=0} = C_4 + \delta B[h_0]$, and the weight $\tfrac12$ there follows
from the first variation of the coefficients of the zero-mode equation,
with no action identity. A weight $w$ on $\mathcal{P}$ would move $K_0$ by
$-1.105\,(w - \tfrac12)\,C_4$ at $\beta = 45$, so the agreement
$K_{G=0} = K_0$ fixes $w = \tfrac12$ to $10^{-8}$. That
agreement is not automatic: at small $G$ the gauge $H_{tt} = H_{xx} = 0$
is reached by a diffeomorphism along $x$ of size $1/G$, which drags the
photon with it and leaves a finite pure-gauge induced field as $s \to 0$.
By the gauge invariance above the pairing does not see it, and since the
boundary metric is held fixed at every $G$, the limit lands on the
fixed-temperature solution.

\subsection{The throat limit}
\label{app:response-throat}

In the fixed-temperature frame, $\beta \to \infty$
takes the interior to $U \to (r - 1)(3r + 1)$, $e^{2Z} \propto (r - \tfrac13)^2$
and $e^{2V} \to \sqrt{\beta/6}$, that is BTZ with $h = 2/3$, $r_0 = 1/3$. With $\rho = r - r_0$, every ingredient of the kernel is governed
by the single radial operator
\begin{equation}
L_m\,y = \big[\rho(\rho^2 - h^2)\,y'\big]' - m\,\rho\,y~,
\end{equation}
which in $\varsigma = h^2/\rho^2$ becomes
$4\varsigma^2(1-\varsigma)y'' - 4\varsigma^2y' - my = 0$, with exponents
$\{0, 0\}$ at the horizon and $(1 \mp \mu)/2$ at the top, $\mu = \sqrt{1+m}$,
and the horizon-regular and normalisable solutions are
$y_H = \varsigma^a\,{}_2F_1(a, a; 1; 1 - \varsigma)$ and $y_I = \varsigma^a\,{}_2F_1(a, a; 1 + \mu; \varsigma)$,
with $a = (1+\mu)/2$. The zero mode $L_{-1}w_0 = 0$ is $m = -1$, $\mu = 0$,
for which the horizon-regular solution is
$w_0 \propto \tfrac{2}{\pi}\sqrt{\varsigma}\,\mathbf{K}(1 - \varsigma)$, with
$\mathbf{K}$ the complete elliptic integral of the first kind; its amplitude
is not normalisable in the throat and cancels from every lattice comparison,
while the quartic measure $p_1 w_0^4\,dr$ converges at the top of the throat.
At finite $\beta$ the same formula
with $a = (1 + i\nu)/2$ supplies the horizon phase $\chi(\nu)$ of section~\ref{sec:4.6}.
The induced field at harmonic $s$ obeys $L_\lambda\hat b = -\lambda\rho
w_0^2$ with $\lambda = s\ell_3^2 e^{-2V} = \gamma b$. The abelian throat
kernel $\mathcal{S}(\gamma) = C_4 - X_{\rm ab}$ is taken in the throat
normalisation of the zero mode, so that $\mathcal{S}(0)$ is the throat value
of $C_4$ and $\mathcal{S}/\mathcal{S}(0)$ is the abelian-only ratio
$(C_4 - X_{\rm ab})/C_4$. It has
the closed-form endpoints $\mathcal{S}(0) = \tfrac12
\int_0^1 \mathcal{F}^4 d\varsigma = 1.89804$, $\mathcal{F} = {}_2F_1(\tfrac12, \tfrac12; 1; 1 - \varsigma)$, and
$\gamma\mathcal{S} \to \tfrac43\mathcal{S}(0)$; its numerical evaluation is
described in appendix~\ref{app:numerics-scan}.

For the metric sector, on the exact throat the solved equations for
$H_{yy}''$ and $\hat b''$ of appendix~\ref{app:response-metric} do not
involve $H_{zz}$ or $H_{zz}'$. With $\Phi = (H_{yy}, \hat b)$ they close on
the block $[\rho(\rho^2 - h^2)\Phi']' = \rho\,(\mathbb{M}\Phi + \mathbb{S})$
with the constant matrix
\begin{equation}
\mathbb{M}_{11} = \mathbb{M}_{22} = \lambda + \tfrac83~, \qquad
\mathbb{M}_{12} = -\frac{16}{9bv}~, \qquad \mathbb{M}_{21} = -b\,(G^2 + 4v)~.
\end{equation}
Since $G^2/v = 3\lambda$ (from the definition of $\lambda$, with
$\ell_3^2 = 1/3$), $\mathbb{M}_{12}\mathbb{M}_{21} = \tfrac{16}{9}(3\lambda + 4) > 0$.
The rescaling $\hat b \to \hat b/\sigma$ with $\sigma = \tfrac34
bv\sqrt{3\lambda + 4}$ then gives the real symmetric form $(\lambda + \tfrac83)
\mathbf{1} - \tfrac43\sqrt{3\lambda + 4}\,\sigma_1$, with
$\lambda$-independent eigenvectors $(1, \mp1)$ and eigenvalues
\begin{equation}
m_\pm(\lambda) = \frac{\big(\sqrt{3\lambda + 4} \pm 2\big)^2}{3}~.
\end{equation}
Neither reaches the Breitenlohner--Freedman bound at any $\lambda$, while
$m_- \le \lambda$ for all $\lambda \ge 0$. The component
$H_{zz}$ is driven by the other fields without feeding back: $z$ lies along
the $\mathrm{AdS}_3$ factor, where gravity has no propagating mode.
The full throat kernel is computed from the complete response system of
appendix~\ref{app:response-metric} on the throat, not from this reduced block, and its
agreement with the brane ladder is documented in appendix~\ref{app:numerics-scan}.

\subsection{The second variation about the uniform lattice}
\label{app:response-hessian}

The stability analysis of section~\ref{sec:5.3} needs the second variation
of the quartic functional about the uniform triangular lattice, over all
lowest-Landau-level perturbations at one flux quantum per cell. Because the
unperturbed density is periodic, the Hessian is block diagonal over the
magnetic Brillouin zone: a perturbation of Bloch momentum $q$ mixes only
with its partner at $-q$, and the $2\times2$ block has eigenvalues, in
units $B = 1$,
\begin{equation}
\begin{split}
&E_\pm(q) = \mathcal{D}(q) + \Sigma_0(q) \pm |\Sigma_2(q)|~, \qquad
\mathcal{D} = \sum_{G\neq0}e^{-G^2/2}K(G^2)\big[\cos(G\wedge q) - 1\big]~,\\
&\Sigma_0 = \sum_G t_G~,\qquad
\Sigma_2 = \sum_G t_G\,e^{-iG\wedge q}~, \qquad
t_G = e^{-|G-q|^2/2}K\big(|G-q|^2\big)~,
\end{split}
\end{equation}
where $G$ runs over the reciprocal lattice and $G \wedge q$ is the
two-dimensional cross product. Stationarity of the uniform amplitude sets
$\epsilon = (K_0 + 2C_\triangle)\eta^2$; its $K_0$ part cancels the
interaction of the perturbation's uniform density with the condensate,
the only place $K_0$ enters, and its $2C_\triangle$ part is the $-1$ in
$\mathcal{D}$. As $q \to 0$, $E_-$ is the Goldstone mode of
translations and rises as $|q|^4$, the soft shear mode of the flat
lowest-Landau-level lattice~\cite{Eilenberger:1967zz,Rosenstein:2010rmp}, while $E_+ \to 2(K_0 + 2C_\triangle)$,
twice the inverse compressibility of section~\ref{sec:5.3}. The spectrum
was evaluated on a dense grid over the zone at fifteen couplings from the
probe limit to $\alpha = 0.385$, twelve of them below $\alpha_L$, and
checked against a direct diagonalisation of the Hessian on finite tori
(table~\ref{tab:checks}).
